\section{Artifact index and reproduction scope}
\label{app:artifacts}

Paths below are relative to the Paper 9 repository root, the parent of
\path{revision-2/}. The complete list of consumed files, byte lengths,
and SHA-256 values is in \path{revision-2/data/source_manifest.json};
the flat index is \path{revision-2/data/source-index.csv}. The manifest
binds 118 inputs to the secondary-analysis script. These bindings identify
the consumed evidence; they are not a replacement for the original
campaign's collection and admission records.

\begin{description}
\item[\hypertarget{artifact-E1}{E1: decoy-gradient channel.}]
\path{evidence/decoy-gradient/paper1-main.tex}, especially the abstract,
the mechanism section, and the mitigation results. The frame counts,
content-probe effect, and utility cost are source-reported from the
published study \citep{paper1}; this revision does not recount all frames.

\item[\hypertarget{artifact-B1}{B1: boundary-replacement audit.}]
\path{evidence/boundary-replacement/independent_evidence_review_2026-09-09/VERDICTS.md}
and its accompanying recomputed-evidence JSON. These records link the
target-alignment defect and utility/control comparisons to their original
runner and outputs. The new analysis hashes these sources but does not
refit corrected attackers or independently repeat all prior audit checks.

\item[\hypertarget{artifact-R1}{R1: trilemma exploration.}]
\path{reproduction/trilemma-campaign/evidence/r1-exploration-20260917/analysis/exploratory_analysis.json},
\path{reproduction/trilemma-campaign/evidence/r1-qualification-20260916/attack-spec-frozen.json}, and
\path{evidence/trilemma-campaign/r1_cell_results.csv} and
\path{evidence/trilemma-campaign/r1_seed_results.csv}.
The script reconstructs 76 cell--seed contrasts, per-binding privacy rules,
utility/gain decisions, and their consistency with the saved tables.
\path{reproduction/trilemma-campaign/evidence/r1-exploration-20260917/confirmation_power.json}
is retained as conditional planning evidence, not an executed confirmation.

\item[\hypertarget{artifact-P4}{P4: defense-redesign pilot.}]
\path{reproduction/defense-redesign-pilot/outputs/runs/pilot-20260922-a/verification.json},
the adjacent report and run-metrics records, and
\path{evidence/defense-redesign-pilot/PILOT_DECISION.md}.
The script recomputes paired utility/gain statistics, candidate advantages,
and recorded qualification/composite decisions. The candidate arrays and
saved qualification scores are the input; raw tensor captures are not
rescored.

\item[\hypertarget{artifact-E8}{E8: orthogonal-wrapper robustness.}]
\path{reproduction/structural-e8/paper-data/collected/h100-2/issues/issue_e8_27b_20260807/e8_robustness_27b_bf16_densek_20260807T230529Z.json}
and \path{evidence/structural-e8/EXPERIMENT_STORY.md}.
The new analysis extracts recorded baselines, depth-specific crossings,
and pooled/session values; it does not generate a new known-pair curve.

\item[\hypertarget{artifact-ER}{ER: deployment and rotation.}]
\path{reproduction/structural-e8/paper-data/collected/toula/er_20260811/COLLECTION_MANIFEST.json},
its admitted summaries and session identity records, and
\path{evidence/structural-e8/EVIDENCE_EXECUTION_LEDGER.md}.
The script extracts eight runs and 24 sessions. The N=16/RTT=0/session-1
identity record contains the reported $2.8$ versus $7.6$ tokens/s rates;
their ratio gives a $63.16\%$ throughput reduction. Recorded precision
does not imply an additional hardware-measurement accuracy guarantee.

\item[\hypertarget{artifact-MI}{MI: estimator and capacity interpretation.}]
\path{evidence/privacy-contract/privacy-contract-draft.tex} and
\path{evidence/privacy-contract/KEY_NUMBERS.md}, with
\path{evidence/structural-e8/LATENT_CAMPAIGN_REPORT_V7_V9.md} for the
finite capacity ablations. The near-zero scores are source-reported.
Equation~\eqref{eq:dv} and the constant-critic counterexample are analytic
arguments; they are not new model measurements.
\end{description}

\section{Executable consistency checks}
\label{app:checks}

From \path{revision-2/}, run
\begin{quote}\small\ttfamily
python3 analysis/recompute.py
\end{quote}
The script uses the Python standard library, writes derived records only
under \path{revision-2/data/}, and loads no checkpoints, tensor captures,
or secret keys. Its stored log reports 1,522 arithmetic and disposition
checks on this input set. They include agreement of independently
reconstructed seed contrasts, confidence endpoints, per-binding rules,
composite outcomes, and the previously generated tables. They are internal
consistency checks, not independent experimental replications or an
estimate of the audit's accuracy on unseen studies.

\begin{table}[h]
\centering\small
\begin{tabular}{@{}lr@{}}
\toprule
R1 rejection-reason intersection & Configurations \\
\midrule
Qualified privacy breach only & 10 \\
Gain failure and qualified privacy breach & 3 \\
Utility and gain failure & 3 \\
Utility failure only & 2 \\
Gain failure only & 1 \\
\midrule
Total & 19 \\
\bottomrule
\end{tabular}
\caption{\textbf{The R1 reason intersections reconstruct the full cohort.}
These disjoint categories explain why marginal reason counts in
Figure~\ref{fig:campaigns} overlap. The six configurations without a
qualified privacy breach have independent utility or gain rejection
reasons. Source: \artifact{R1}; generated \texttt{data/audit\_results.json}.}
\label{tab:intersections}
\end{table}

The derived JSON also preserves uncertainty conventions, denominators,
source levels, and unresolved historical limitations. The figure script
checks the generated CSVs against the original cell/seed tables and pilot
report before writing vector PDFs. Regeneration commands and dependencies
are documented in the revision README. An additional preservation manifest
checks the existing \path{paper/} files independently of the new revision.
